%%
%% * DOCUMENT:
%%   - title: `PhD report on Interactive Clustering`
%%   - autor: Erwan SCHILD
%%   - date: 02/09/2022
%% * LATEX TEMPLATE:
%%   - title: thesul v0.15
%%   - autor: Denis ROEGEL
%%   - date: 30 mars 2013, update 12 mars 2022
%%   - url: [thesul](https://members.loria.fr/DRoegel/TeX/TUL.html)
%%

%%%%%%%%%%%%%%%%%%%%%%%%%%%%%%%%%%%%%%%%%%%%%%%%%%%%%%%%%%%%%%%%%%%%%%%
%%%%% PARAMÈTRES DU DOCUMENT
%%%%%%%%%%%%%%%%%%%%%%%%%%%%%%%%%%%%%%%%%%%%%%%%%%%%%%%%%%%%%%%%%%%%%%%

%%% Pour la classe du document:
\documentclass[11pt]{thesul}

%%% Pour déclarer un brouillon:
% \ThesisDraft

%%% Pour les vraies marges:
% \SetRealMargins{1mm}{1mm}


%%%%%--------------------------------------------------------------------
%%%%% Chargement des packages
%%%%%--------------------------------------------------------------------

%%% Pour un PDF avec des hyperliens:
\usepackage[pageanchor=false]{tulhypref}

%%% Pour une bibliographie au styme 'named':
% \usepackage{named}

%%% Pour les figures:
\usepackage{graphicx}

%%% Pour les mini-tables des matières:
\usepackage[french]{minitoc}


%%%%%--------------------------------------------------------------------
%%%%% Paramètres des en-têtes
%%%%%--------------------------------------------------------------------

%%% Exemples d'en-têtes:
% \UppercaseHeadings 
% \UnderlineHeadings
% \newcommand\bfheadings[1]{{\bf #1}}
% \FormatHeadingsWith{\bfheadings}
% \FormatHeadingsWith{\uppercase}
% \FormatHeadingsWith{\underline}

\newcommand\upun[1]{\uppercase{\underline{\underline{#1}}}}
\FormatHeadingsWith\upun
\newcommand\itheadings[1]{\textit{#1}}
\FormatHeadingsWith{\itheadings}

% Pour un trait sous une en-tete:
\setlength{\HeadRuleWidth}{0.4pt}


%%%%%--------------------------------------------------------------------
%%%%% Paramètres des référence, du glossaire et de l'index
%%%%%--------------------------------------------------------------------

%%% Enlever les numéro et préfix des chapitres:
% \NoChapterNumberInRef
% \NoChapterPrefix

%%% Faire un glossaire:
\makeglossary

%%% Faire un index:
\makeindex


%%%%%%%%%%%%%%%%%%%%%%%%%%%%%%%%%%%%%%%%%%%%%%%%%%%%%%%%%%%%%%%%%%%%%%%
%%%%% DÉBUT DU DOCUMENT
%%%%%%%%%%%%%%%%%%%%%%%%%%%%%%%%%%%%%%%%%%%%%%%%%%%%%%%%%%%%%%%%%%%%%%%

\begin{document}

%%% Pour définir le style de pages:
\FramePartsInToc
\OddHead={{\leftmark\rightmark}{\hfil\slshape\rightmark}}
\EvenHead={{\leftmark}{{\slshape\leftmark}\hfil}}
\OddFoot={\hfil\thepage}
\EvenFoot={\thepage\hfil}
\pagestyle{ThesisHeadingsII}

%%%%%--------------------------------------------------------------------
%%%%% Encadrement des partie set des chapitres
%%%%%--------------------------------------------------------------------

%%% Pour encadrer encadre les parties dans la table des matières:
%%% (nb: ces commandes doivent figurer apres \begin{document})
\FramePartsInToc

%%% Pour encadrer encadre les chapitres dans la table des matières:
%%% (nb: ces commandes doivent figurer apres \begin{document})
% \FrameChaptersInToc

%-------------------------------------------------------------------
%            Réinitialisation de la numérotation des chapitres
%-------------------------------------------------------------------

% Si la commande suivante est présente, elle doit figurer APRÈS '\begin{document}' et avant la première commande '\part':
% \ResetChaptersAtParts

%-------------------------------------------------------------------
%               mini-tables des matières par chapitre
%-------------------------------------------------------------------

% Préparer les mini-tables des matières par chapitre (commande de minitoc.sty):
\dominitoc

%%%%%--------------------------------------------------------------------
%%%%% Page de titre
%%%%%--------------------------------------------------------------------

%%% Titre:
\ThesisTitle{
    Faciliter la conception d'un assistant conversationnel avec le clustering interactif
}

%%% Date:
\ThesisDate{09 setpembre 2022}

%%% Auteur:
\ThesisAuthor{Erwan SCHILD}

%%% Type de these:
\ThesisUL

%%% Membres du jury:
\President={
M. Le président de la soutenance
}
\Rapporteurs ={
Le rapporteur 1\\
Le rapporteur 2
}
\Examinateurs={
L'examinateur 1\\
L'examinateur 2
}
\Encadrants={
Dr. Jean-Charles LAMIREL\\
Dr. Florian MICONI
}

%%% Création de la page de titre:
\MakeThesisTitlePage

%%%%%--------------------------------------------------------------------
%%%%% Page de remerciements
%%%%%--------------------------------------------------------------------

\begin{ThesisAcknowledgments}
    Les remerciements.\\
    Ma femme.\\
    Mes amis.\\
    Mon chien.\\
\end{ThesisAcknowledgments}

%%%%%--------------------------------------------------------------------
%%%%% Page de dédicaces
%%%%%--------------------------------------------------------------------

\begin{ThesisDedication}
    Je dédie cette thèse\\
    à ma machine.\\
    Oui, à Pandore,\\
    qui fut la première de toutes.
\end{ThesisDedication}

%%%%%--------------------------------------------------------------------
%%%%% Table des matières
%%%%%--------------------------------------------------------------------

%%% Ajout des parties dans la table des matières:
\WritePartLabelInToc

%%% Ajout des chapitres dans la table des matières:
\WriteChapterLabelInToc

%%% Génération de la table des matières:
\tableofcontents

%%% Génère la liste des figures:
\listoffigures

%%%%%%%%%%%%%%%%%%%%%%%%%%%%%%%%%%%%%%%%%%%%%%%%%%%%%%%%%%%%%%%%%%%%%%%
%%%%% PRÉAMBULE
%%%%%%%%%%%%%%%%%%%%%%%%%%%%%%%%%%%%%%%%%%%%%%%%%%%%%%%%%%%%%%%%%%%%%%%

\SpecialSection{Avertissement}

Attention à ce que vous lisez !

\SpecialSection{Préambule}

Ceci est un ensemble d'information complémentaires à savoir...

%%% Commencer la numérotation arabe (cf. '\pagenumbering{arabic}') avec la page 1 sur une page impaire:
\mainmatter

%%%%%%%%%%%%%%%%%%%%%%%%%%%%%%%%%%%%%%%%%%%%%%%%%%%%%%%%%%%%%%%%%%%%%%%
%%%%% PARTIE 1
%%%%%%%%%%%%%%%%%%%%%%%%%%%%%%%%%%%%%%%%%%%%%%%%%%%%%%%%%%%%%%%%%%%%%%%

\chapter{Introduction}
    \label{chapter:1-INTRODUCTION}
	\todo[inline]{CHAPITRE À REFORMULER FAÇON SWALES}

    %%%%%--------------------------------------------------------------------
    %%%%% Section 1.1:
    %%%%%--------------------------------------------------------------------
    \section{"\textit{Asset centrality}"}
	\todo[inline]{SECTION À RÉDIGER}

    \begin{todolist}
    
        \item Des enjeux ou problèmes actuels

            •	Accessibilité à l'information :
                o	Grosses bases documentaires, pas toujours ordonnées ;

            •	Relations client à distance
                o   Besoin d'un accessibilité h24 ;

        \item Utilisation de plus en plus fréquente des chatbots
        
            •   Description succinte ;

            •   Cas d'usage usuels ;

            •   Tous les canaux d'utilisation ;

            •   Avantages et Dérives potentiels de l'utilisation (emploi, biais, pertinence, ergonomie, ...) ;

        \item Révolution techniques fréquentes (règles, classification, modèles)

            •	Moteurs de règles :
                o	Basé sur la détecté de mots clés,
                o	(+) facile à mettre en œuvre,
                o	(-) peu robuste au langage naturel,
                o	Paramétrage des réponses ;

            •	Paramétrage intentions-entités :
                o	Classification d’intention et/ou détection d’entités,
                o	(+) plus robuste au langage naturel, facile à paramétrer, réponses contrôlées,
                o	(-) demande de l’entrainement, des données, …,
                o	Paramétrage des réponses ;

            •	Génération de réponse :
                o	Réseau de neurones avec attention,
                o	Transformers,
                o	(+) plus robuste,
                o	(-) plus complexe à mettre en œuvre, réponses non contrôlées,
                o	Réponses non paramétrées ;

            •	Approche hybride :
                o   Cumul des trois approches pour cumuler certains avantages suivant les besoins ;
	\end{todolist}

    %%%%%--------------------------------------------------------------------
    %%%%% Section 1.2:
    %%%%%--------------------------------------------------------------------
	
	\section{"\textit{Estabilishing a Niche}"}
	\todo[inline]{SECTION À RÉDIGER}

    \begin{todolist}

        \item Cadre industriel

            •	Algorithme fixe

            •	Données spécifiques

        \item GAP: Besoins de données

            •	Collecte de données spécifiqus au domaine traité :
                o   extraction de base de données (solution simple),
                o   collecte manuelle (organisation complexe, biais de collecte),
                o   scraping (pas toujours fiable) ;
               \todo[inline]{
	               Remarque Gautier 20/02/2023: utilité du travail
	               Un aspect à réfléchir ici : on a besoin de données, en effet, et par conséquent on génère une industrie de l'annotation. Tout se passe un peu comme si on déportait tout le travail nécessaire pour accompagner les clients qui utilisent le chatbot sur les phases d'nanotation. Ca pose une question importante de l'utilité du travail : travaille-t-on pour l'humain ou pour la machine ? (ca permet d'aborder la question des débats anti-IA aussi)
	               Pour éviter la déshumanisation du travail, c'est donc très important de réduire l'adhérence aux données et le besoin d'annotation.
               }

            •	Nombreux biais :
                o   Biais,,
                o   Réglementation,
                o   Compétences (NOTRE COEUR DU SUJET),
                o   ...
    \end{todolist}

    %%%%%--------------------------------------------------------------------
    %%%%% Section 1.3:
    %%%%%--------------------------------------------------------------------
	
	\section{"\textit{Occupying the Niche}"}
	\todo[inline]{SECTION À RÉDIGER}

    \begin{todolist}
        \item Etude de l'organisation d'une entreprise pour concevoir ses jeux de données
        \item Etude de l'état de l'art pour concevoir des jeux de données
        \item proposition/contribution : une méthode adaptée pour un cadre industriel
    \end{todolist}

%%%%%%%%%%%%%%%%%%%%%%%%%%%%%%%%%%%%%%%%%%%%%%%%%%%%%%%%%%%%%%%%%%%%%%%
%%%%% PARTIE 2
%%%%%%%%%%%%%%%%%%%%%%%%%%%%%%%%%%%%%%%%%%%%%%%%%%%%%%%%%%%%%%%%%%%%%%%

\part{État de l'art (La seconde partie)}
    \label{PARTIE_2_ETAT_DE_L_ART}
\minitoc

\chapter{Dans la fiction}

    \section{Rêvons un peu...}
        Une « autre » page avec « plein » de texte « et » très varié.
        Une autre page avec plein de texte très varié .
        Une autre page avec plein de texte très varié.

    \section{Dans les films...}
        Une « autre » page avec « plein » de texte « et » très varié.
        Une autre page avec plein de texte très varié .
        Une autre page avec plein de texte très varié.

    \section{Dans les comics...}
        Une « autre » page avec « plein » de texte « et » très varié.
        Une autre page avec plein de texte très varié .
        Une autre page avec plein de texte très varié.

\chapter{Dans la recherche}

    \section{Depuis 1980...}
        Une autre page avec plein de texte très varié.
        Une autre page avec plein de texte très varié.

    \section{Depuis 1990...}
        Une autre page avec plein de texte très varié.
        Une autre page avec plein de texte très varié.

    \section{Depuis 2000...}
        Une autre page avec plein de texte très varié.
        Une autre page avec plein de texte très varié.

    \section{Depuis 2010...}
        Une autre page avec plein de texte très varié.
        Une autre page avec plein de texte très varié.

        \begin{figure}[htb]
            \includegraphics{tulul}
            \caption{x\label{rintintin}}
        \end{figure}

    \section{Depuis 2022...}
        Une autre page avec plein de texte très varié.
        Une autre page avec plein de texte très varié.

\chapter{Dans l'industrie}

    \section{Avec XXXXX...}
        Une autre page avec plein de texte très varié.
        Une autre page avec plein de texte très varié.

    \section{Avec YYYY...}
        Une autre page avec plein de texte très varié.
        Une autre page avec plein de texte très varié.

    \section{Depuis ZZZZZ...}

        \begin{itemize}
            \item Une autre page avec plein de texte très varié.
            \item Une autre page avec plein de texte très varié.
                \begin{itemize}
                    \item Une autre page.
                    \item Une autre page.
                \end{itemize}
            \item Une autre page avec plein de texte très varié.
        \end{itemize}

%%%%%%%%%%%%%%%%%%%%%%%%%%%%%%%%%%%%%%%%%%%%%%%%%%%%%%%%%%%%%%%%%%%%%%%
%%%%% ANNEXES
%%%%%%%%%%%%%%%%%%%%%%%%%%%%%%%%%%%%%%%%%%%%%%%%%%%%%%%%%%%%%%%%%%%%%%%

\FrameThisInToc
\Annexes
    \label{ANNEXES}
    \Annexes

\Annex{Annexe théorique}
\label{annex:A-ANNEXE-THEORIQUE}

	\minitoc

    \section{Les algorithmes de clustering }

        \subsection{Kmeans}
        kmeans \index{clustering!kmeans}

        \subsection{Hierarchique}
        hierarchique \index{clustering!hierarchique}

        \subsection{Spectral}
        spectral \index{clustering!spectral}

        \subsection{DBScan}
        dbscan \index{clustering!dbscan}

        \subsection{Affinity Propagation}
        affinity propagation \index{clustering!affinity propagation}


    \section{Evaluation d’une clustering}

        \subsection{Homogénéité – Complétude – Vmeasure}
        la VMeasure \index{vmeasure} est la moyenne harmonique entre l'homogénéité et la complétude.
        
        \subsection{FMC}

\Annex{Annexe technique}
\label{annex:B-ANNEXE-TECHNIQUE}

	\minitoc

    \section{package pypi interactive-clustering}
    \section{package pypi interactive-clustering-gui}
    \section{package pypi features-maximization-metrics}
    \section{experimentations jupyter notebook}

\Annex{Annexe des jeux de données}
\label{annex:C-ANNEXE-DATASET}

	\minitoc

    \section{BANK CARDS: french bank cards}
	\label{annex:C.1-DATASET-BANK-CARDS}
	
		[FR] Jeu d'entraînement en français d'assistants conversationnels traitant des demandes courantes sur les cartes bancaires.
		
		Description : Cet ensemble de données représente des exemples de demandes usuelles des clients concernant la gestion des cartes bancaires. Il peut être utilisé comme jeu d'entraînement pour un assistant conversationnel destiné à traiter ces demandes courantes.
		
		Contenu : Les questions sont formulées en français. L'ensemble de données est divisé en 10 intentions de 50 questions chacune, pour un total de 500 questions.
		
		Périmètre des intentions : Les intentions sont construites de telle manière que toutes les questions issues d'une même intention ont la même réponse ou action. Le périmètre couvert concerne : la perte ou le vol de cartes ; la carte avalée ; la commande des cartes ; la consultation du solde bancaire ; l'assurance fournie par une carte ; le déverrouillage de la carte ; la gestion de cartes virtuelles ; la gestion du découvert bancaire ; la gestion des plafonds de paiement ; la gestion du mode sans contact.
		
		Origine : Le périmètre des intentions est inspiré par un chatbot actuellement en production, et la formulation des questions est inspirée de demandes courantes de clients.
	
    \section{MLSUM: press titles}
	\label{annex:C.2-DATASET-MLSUM-SUBSET-SCHILD}
	
		We present MLSUM, the first large-scale MultiLingual SUMmarization dataset.
		Obtained from online newspapers, it contains 1.5M+ article/summary pairs in five different languages -- namely, French, German, Spanish, Russian, Turkish. Together with English newspapers from the popular CNN/Daily mail dataset, the collected data form a large scale multilingual dataset which can enable new research directions for the text summarization community. We report cross-lingual comparative analyses based on state-of-the-art systems. These highlight existing biases which motivate the use of a multi-lingual dataset.
		
		For constraints annotation experiment based on data similarity, this dataset have been subsetted (randomly pick 75 articles in the following 14 most used topics: 'economie', 'politique', 'sport', 'planete' (renamed in 'ecologie'), 'sciences', 'police-justice', 'disparitions', 'emploi', 'sante', 'musiques', 'arts', 'educations', 'climat' (renamed in 'meteo'), 'immobilier') and filtered (keep articles that have an obvious topics regarding their titles, without their bodies). Two reviewers have working on this task in order to limit the subjectivity of the filtering. This subsetted dataset is used (1) to estimate needed time to annotate titles similarity with constraints (MUST-LINK, CANNOT-LINK) and (2) to test interactive clustering methodology (constraints annotation and constrained clustering).


    %%%%%%%%%%%%%%%%%%%%%%%%%%%%%%%%%%%%%%%%%%%%%%%%%%%%%%%%%%%%%%%%%%%%%%%
    %%%%% GLOSSAIRE
    %%%%%%%%%%%%%%%%%%%%%%%%%%%%%%%%%%%%%%%%%%%%%%%%%%%%%%%%%%%%%%%%%%%%%%%

    %%% Définition du glossaire:
    \Glossary{Chat1}{animal}\Glossary{Chien1}{Autre animal}
    \Glossary{Chat2}{animal}\Glossary{Chien2}{Autre animal}
    \Glossary{Chat3}{animal}\Glossary{Chien3}{Autre animal}
    \Glossary{Chat4}{animal}\Glossary{Chien4}{Autre animal}
    \Glossary{Chat5}{animal}\Glossary{Chien5}{Autre animal}
    \Glossary{Chat6}{animal}\Glossary{Chien6}{Autre animal}
    \Glossary{Chat7}{animal}\Glossary{Chien7}{Autre animal}
    \Glossary{Chat8}{animal}\Glossary{Chien8}{Autre animal}
    \Glossary{Chat9}{animal}\Glossary{Chien9}{Autre animal}
    \Glossary{Chat10}{animal}\Glossary{Chien10}{Autre animal}
    \Glossary{Chat11}{animal}\Glossary{Chien11}{Autre animal}
    \Glossary{Chat12}{animal}\Glossary{Chien12}{Autre animal}
    \Glossary{Chat13}{animal}\Glossary{Chien13}{Autre animal}
    \Glossary{Chat14}{animal}\Glossary{Chien14}{Autre animal}
    \Glossary{Chat15}{animal}\Glossary{Chien15}{Autre animal}
    \Glossary{Chat16}{animal}\Glossary{Chien16}{Autre animal}
    \Glossary{Chat17}{animal}\Glossary{Chien17}{Autre animal}
    \Glossary{Chat18}{animal}\Glossary{Chien18}{Autre animal}
    \Glossary{Chat19}{animal}\Glossary{Chien19}{Autre animal}
    \Glossary{Chat22}{animal}\Glossary{Chien22}{Autre animal}
    \Glossary{Chat23}{animal}\Glossary{Chien23}{Autre animal}
    \Glossary{Chat24}{animal}\Glossary{Chien24}{Autre animal}
    \Glossary{Chat25}{animal}\Glossary{Chien25}{Autre animal}
    \Glossary{Chat26}{animal}\Glossary{Chien26}{Autre animal}
    \Glossary{Chat27}{animal}\Glossary{Chien27}{Autre animal}
    \Glossary{Chat28}{animal}\Glossary{Chien28}{Autre animal}
    \Glossary{Chat29}{animal}\Glossary{Chien29}{Autre animal}

    %%% Affichage du glossaire:
    \BeginGloWith{
        Voici un glossaire tout-à-fait fictif,
        introduit par un texte sur toute la largeur
        des deux colonnes.
    }
    \twocolumn
    \PrintGlossary

    %%%%%%%%%%%%%%%%%%%%%%%%%%%%%%%%%%%%%%%%%%%%%%%%%%%%%%%%%%%%%%%%%%%%%%%
    %%%%% INDEX
    %%%%%%%%%%%%%%%%%%%%%%%%%%%%%%%%%%%%%%%%%%%%%%%%%%%%%%%%%%%%%%%%%%%%%%%

    %%% Affichage de l'index:
    \BeginIndWith{
        Voici un index
    }
    \PrintIndex

    \onecolumn

    %%%%%%%%%%%%%%%%%%%%%%%%%%%%%%%%%%%%%%%%%%%%%%%%%%%%%%%%%%%%%%%%%%%%%%%
    %%%%% BIBLIOGRAPHIE
    %%%%%%%%%%%%%%%%%%%%%%%%%%%%%%%%%%%%%%%%%%%%%%%%%%%%%%%%%%%%%%%%%%%%%%%

    %%% Définition de la bibliographie:
    \begin{thebibliography}{Lam91a}

        \bibitem[CM88]{chandy88a}
        K.~M. Chandy and J.~Misra.
        \newblock \emph{Parallel program design: a foundation}.
        \newblock Addison-Wesley Publishing Company, 1988.

        \bibitem[Lam91a]{lamport91a}
        Leslie Lamport.
        \newblock {The Temporal Logic of Actions}.
        \newblock Technical Report~79, SRC, 1991.

    \end{thebibliography}

    %%%%%%%%%%%%%%%%%%%%%%%%%%%%%%%%%%%%%%%%%%%%%%%%%%%%%%%%%%%%%%%%%%%%%%%
    %%%%% ABSTRACT
    %%%%%%%%%%%%%%%%%%%%%%%%%%%%%%%%%%%%%%%%%%%%%%%%%%%%%%%%%%%%%%%%%%%%%%%

    % (si le résumé apparaît sur une colonne étroite, avec la bibliographie à gauche, c'est sans doute parce que vous avez oublié de générer les fichiers d'index et de glossaire...)

    \NumberAbstractPages
    \begin{ThesisAbstract}

        %%%%%--------------------------------------------------------------------
        %%%%% Abstract en français
        %%%%%--------------------------------------------------------------------
        \begin{FrenchAbstract}
            Le résumé.
            \KeyWords{chat, chien, puces.}
        \end{FrenchAbstract}

        %%%%%--------------------------------------------------------------------
        %%%%% Abstract en anglais
        %%%%%--------------------------------------------------------------------
        \begin{EnglishAbstract}
            In computational geometry many search problems and range queries can be solved by performing an iterative search for the same key in separate ordered lists.  In Part I of this report we show that, if these ordered lists can be put in a one-to-one correspondence with the nodes of a graph of degree  d  so that the iterative search always proceeds along edges of that graph, then we can do much better than the obvious sequence of binary searches. Without expanding the storage by more than a constant factor, we can build a data-structure, called a fractional cascading structure, in which all original searches after the first can be carried out at only  log d  extra cost per search.  Several results related to the dynamization of this structure are also presented. Part II gives  numerous applications of this technique to geometric problems.
            Examples include intersecting a polygonal path with a line, slanted range search, orthogonal range search, computing locus functions, and others. Some results on the optimality of fractional cascading, and certain extensions of the technique for retrieving additional information are also included.
            \KeyWords{cat, dog, flees.}
        \end{EnglishAbstract}
    \end{ThesisAbstract}

\end{document}